\documentclass[11pt]{article}

\usepackage[margin=1in]{geometry}
\usepackage{booktabs}
\usepackage{hyperref}
\usepackage{longtable}
\usepackage{tabularx}
\usepackage{xcolor}

\hypersetup{
  colorlinks=true,
  linkcolor=blue,
  urlcolor=blue
}

\title{Technical Project Overview\\
\large LeGrad-Inspired Temporal Explanations for Speech Emotion Recognition}
\author{gradient\_based\_speach\_xai}
\date{June 7, 2026}

\begin{document}
\maketitle

\section{Project Overview}

This project compares two approaches for explaining a pretrained speech emotion recognition classifier. SpeechXAI provides human-readable explanations based on word-aligned audio segments. The proposed LeGrad-inspired method will derive temporal relevance from gradients associated with the attention mechanisms of a wav2vec 2.0 encoder. The classification task uses speech samples and emotion annotations from IEMOCAP.

The quantitative evaluation uses deletion by silence masking. Important audio regions are replaced by silence while the original signal length is preserved, and the resulting confidence drop for the classifier's original predicted class is measured. The central methodological decision is to compare word-level SpeechXAI explanations with time-bin or token-level LeGrad explanations under the same total occluded duration.

\section{Sources and External References}

\begin{itemize}
  \item \textbf{Dataset: IEMOCAP.} The local data follows the five-session IEMOCAP release from the USC Signal Analysis and Interpretation Laboratory. The corpus contains speech, transcripts, and categorical emotion annotations. Project source: \href{https://sail.usc.edu/iemocap/}{USC IEMOCAP website}. Dataset paper: \href{https://sail.usc.edu/iemocap/Busso_2008_iemocap.pdf}{Busso et al., ``IEMOCAP: Interactive Emotional Dyadic Motion Capture''}.

  \item \textbf{Pretrained classifier.} The project uses \href{https://huggingface.co/superb/wav2vec2-base-superb-er}{\texttt{superb/wav2vec2-base-superb-er}}, the Hugging Face wav2vec 2.0 emotion-recognition checkpoint loaded by the SpeechXAI repository for IEMOCAP. The model exposes four output classes: anger, happiness/excitement, neutral, and sadness.

  \item \textbf{wav2vec 2.0.} The encoder architecture originates from Baevski et al., \href{https://arxiv.org/abs/2006.11477}{``wav2vec 2.0: A Framework for Self-Supervised Learning of Speech Representations''}.

  \item \textbf{SUPERB.} The active checkpoint follows the SUPERB speech-processing benchmark described by Yang et al., \href{https://arxiv.org/abs/2105.01051}{``SUPERB: Speech Processing Universal PERformance Benchmark''}.

  \item \textbf{SpeechXAI.} The word-level explanation method follows Pastor et al., \href{https://aclanthology.org/2024.eacl-long.136/}{``Explaining Speech Classification Models via Word-Level Audio Segments and Paralinguistic Features,'' EACL 2024}. Repository: \href{https://github.com/elianap/SpeechXAI}{elianap/SpeechXAI}.

  \item \textbf{WhisperX alignment.} SpeechXAI obtains word timestamps with Bain et al., \href{https://www.isca-archive.org/interspeech_2023/bain23_interspeech.html}{``WhisperX: Time-Accurate Speech Transcription of Long-Form Audio,'' Interspeech 2023}. WhisperX combines VAD-guided segmentation, Whisper transcription, and external phoneme alignment.

  \item \textbf{LeGrad.} The proposed temporal method is inspired by Bousselham et al., \href{https://arxiv.org/abs/2404.03214}{``LeGrad: An Explainability Method for Vision Transformers via Feature Formation Sensitivity''}. Repository: \href{https://github.com/WalBouss/LeGrad}{WalBouss/LeGrad}.
\end{itemize}

The current project configuration uses IEMOCAP. No separate YAMLCap dataset source is configured.

\section{Model Architecture Used by the Current Pipeline}

The executable transformer-relevance scripts in the current branch load the
SUPERB HuBERT emotion-recognition checkpoint through
\texttt{load\_hubert\_emotion\_model}:

\begin{verbatim}
superb/hubert-base-superb-er
\end{verbatim}

The older project configuration still contains references to the wav2vec 2.0
checkpoint used in the SpeechXAI comparison path, but the implemented rollout
and head-relevance pipeline is written around a HuBERT sequence-classification
model. Notebook \path{notebooks/00_model_used.ipynb} documents this architecture
and provides executable inspection cells.

The executed notebook reports the following active configuration values:

\begin{itemize}
  \item model type: \texttt{hubert}
  \item architecture: \texttt{HubertForSequenceClassification}
  \item Transformer depth: 12 layers and 12 attention heads
  \item hidden size: 768
  \item classifier projection size: 256
  \item weighted layer aggregation: \texttt{use\_weighted\_layer\_sum=True}
  \item output labels: \texttt{neu}, \texttt{hap}, \texttt{ang}, and \texttt{sad}
\end{itemize}

The notebook also shows that the older YAML configuration still points to the
following wav2vec 2.0 checkpoint, while the current transformer-relevance
scripts use the HuBERT checkpoint above:

\begin{verbatim}
superb/wav2vec2-base-superb-er
\end{verbatim}

The HuBERT model first converts the waveform into a lower-rate sequence of
acoustic frames with a convolutional feature encoder. These frames are then
processed by the HuBERT Transformer encoder, which produces temporal hidden
states rather than image patches. The convolutional frontend is important for
mapping token relevance back to approximate waveform intervals, but it is not
the main source of difficulty for adapting LeGrad.

The more important issue is the sequence-classification head. In the Hugging
Face HuBERT sequence-classification architecture, the classifier path can be
written schematically as follows:

\begin{verbatim}
waveform
  -> convolutional feature encoder
  -> HuBERT Transformer hidden states
  -> optional learned weighted sum of hidden layers
  -> projector
  -> temporal mean pooling
  -> classifier
  -> emotion logits
\end{verbatim}

When \texttt{use\_weighted\_layer\_sum} is enabled, the head does not consume
one individual Transformer layer directly. Instead, it forms a learned mixture
of hidden layers,
\[
H_{\mathrm{mix}} = \sum_{l=0}^{12} \mathrm{softmax}(\alpha)_l H_l,
\]
and passes this mixed representation into the projection, pooling, and classifier stack. This
matters for LeGrad because the original method defines a target score at each
Transformer layer before differentiating that score with respect to the
attention map of the same layer.

More explicitly, the HuBERT forward pass returns a tuple of hidden-state
tensors:

\begin{verbatim}
outputs.hidden_states = [H_0, H_1, ..., H_12]
\end{verbatim}

Here $H_0$ is the temporal representation entering the Transformer stack after
the convolutional/frontend and feature-projection stages, while
$H_1,\ldots,H_{12}$ are the outputs of the 12 Transformer layers. Each hidden
state has shape
\[
\mathrm{shape}(H_l) = B \times T \times 768,
\]
where $B$ is the batch size, $T$ is the number of HuBERT temporal frames, and
768 is the hidden dimension. The model stacks these tensors along a new layer
axis:

\begin{verbatim}
stacked_hidden_states = stack([H_0, H_1, ..., H_12])
                      # shape: [B, 13, T, 768]
\end{verbatim}

The checkpoint contains one learned scalar per entry in this stack:
\begin{verbatim}
layer_weights  # shape: [13]
\end{verbatim}

These raw scalars are the stored trainable parameters. During each forward pass,
the model normalizes them with a softmax and uses the resulting coefficients to
form the representation that the classification head actually receives:

\begin{verbatim}
norm_weights = softmax(layer_weights)       # shape: [13]

H_mix = sum_l norm_weights[l] * H_l         # shape: [B, T, 768]

projected = projector(H_mix)                # shape: [B, T, 256]
pooled = mean_over_time(projected)           # shape: [B, 256]
logits = classifier(pooled)                 # shape: [B, 4]
\end{verbatim}

The softmax coefficients are therefore not stored as independent checkpoint
parameters. They are recomputed from \texttt{layer\_weights} every time the
model runs, which preserves the same mixture logic if the model is fine-tuned:
fine-tuning changes the raw \texttt{layer\_weights}, and the forward pass then
produces updated normalized coefficients.

For explainability, it is important not to confuse these coefficients with a
LeGrad layer score. The coefficient
$\alpha_l=\mathrm{softmax}(\texttt{layer\_weights})_l$ says how much the
trained HuBERT head mixes representation $H_l$ into $H_{\mathrm{mix}}$. It is a
learned, forward-time mixture coefficient. It is not by itself a
class-dependent score for the current utterance. LeGrad, however, requires a
scalar activation whose gradient can be taken with respect to an attention map,
for example
\[
s_l = \mathrm{classifier}(\mathrm{pool}(\mathrm{projector}(H_l)))_c
\]
or the real final score
\[
s = \mathrm{classifier}(\mathrm{pool}(\mathrm{projector}(H_{\mathrm{mix}})))_c.
\]
If one used only $\alpha_l$ as the score, its derivative with respect to the
attention map $A_l$ would not provide an utterance- or class-specific temporal
explanation, because $\alpha_l$ is not produced from the current attention map.
The learned coefficients are therefore better viewed as possible weights for
aggregating already-computed layer relevance maps, not as the relevance score
from which those maps are derived.

Thus the learned mixture occurs before the projector and before temporal
pooling. The temporal structure is still present in $H_{\mathrm{mix}}$ because
the tensor keeps the $T$ frame axis; the utterance-level compression happens
after projection, when the model averages over time.

The \texttt{projector} should therefore not be confused with the layer-mixing
operation. The layer mixture uses \texttt{layer\_weights} to combine
$H_0,\ldots,H_{12}$ into $H_{\mathrm{mix}}$. The projector is a separate learned
linear layer, with \texttt{projector.weight} shape \texttt{(256, 768)}, applied
to each temporal frame of $H_{\mathrm{mix}}$ to map 768-dimensional HuBERT
features into the 256-dimensional space consumed by the final classifier.

The executed checkpoint-tensor inspection in Notebook
\path{notebooks/00_model_used.ipynb} confirms that
\texttt{layer\_weights} has shape \texttt{(13,)}, while the classifier head has
\texttt{projector.weight} shape \texttt{(256, 768)} and
\texttt{classifier.weight} shape \texttt{(4, 256)}. The 13 mixture entries
correspond to $H_0$ plus the 12 Transformer layer outputs. This is important:
the final HuBERT emotion head mixes representations from all these levels
before classification, whereas attention maps are only available for the 12
Transformer layers. The forward-time softmax coefficients computed from the
inspected raw checkpoint weights are spread across layers rather than
concentrated in one layer; they range from 0.067804 for $H_1$ to 0.084699 for
$H_{10}$.

The notebook's executed source-code inspection confirms this exact forward
path in the installed \texttt{transformers} implementation: the model sets
\texttt{output\_hidden\_states=True} when weighted layer aggregation is active,
stacks all hidden states, applies \texttt{softmax(self.layer\_weights)}, sums
the weighted layers, then applies \texttt{self.projector}, temporal pooling, and
\texttt{self.classifier}. Thus the ordinary classifier score is produced from
the learned layer mixture, not from an isolated intermediate layer.

In the ViT/CLIP setting used by LeGrad, the layer-wise score is more direct.
Each Transformer block produces visual patch tokens in the same representation
stream. The official LeGrad implementation takes an intermediate block output,
aggregates the tokens into an image feature, applies the final visual
normalization/projection, and computes the target image-text score. Thus the
intermediate score follows a natural visual-token pathway:

\begin{verbatim}
ViT layer H_l
  -> token aggregation
  -> final visual projection
  -> target score_l
\end{verbatim}

For HuBERT, applying the classifier head to one individual layer,

\begin{verbatim}
HuBERT layer H_l
  -> projector
  -> temporal mean pooling
  -> classifier
  -> target score_l
\end{verbatim}

is a methodological adaptation. It may be useful and testable, but it may not
exactly match the representation pathway learned by a checkpoint whose head was
trained on $H_{\mathrm{mix}}$ rather than on each single $H_l$. This is the
specific architectural reason why a faithful temporal LeGrad implementation for
HuBERT is less trivial than applying the original ViT code unchanged.

\section{Implemented Methods and Core Logic}

The current branch implements a lightweight temporal-relevance pipeline for a HuBERT-based speech emotion classifier. Rather than attempting full Transformer-aware relevance propagation, the implementation uses a compact approximation that combines two ideas: attention flow through the encoder and class-sensitive evidence from the final classifier head.

After checking the official LeGrad implementation, the terminology should be
kept precise. Original LeGrad uses the positive gradient of the target
activation with respect to an attention map as the explanation signal,
$\max(0, \partial s / \partial A)$. It does not multiply that gradient by the
attention values, and it does not perform attention rollout. The implementation
in this repository should therefore be described as a LeGrad-inspired
attention-gradient rollout variant, not as a direct implementation of LeGrad.

The pipeline has three conceptual stages.

\begin{enumerate}
  \item \textbf{Attention-gradient rollout.} The model is run once with
  attention outputs enabled and gradients are taken with respect to the target
  class logit. For each Transformer layer, positive gradient-weighted attention
  is averaged over heads and then propagated through the layer stack to form a
  joint rollout matrix describing how information can diffuse across tokens.

  \item \textbf{Head/pooling relevance.} The final HuBERT token representations that enter the classification head are scored for a target class using the classifier weight. Under mean pooling, each token receives a contribution proportional to $H_t \cdot w_c / T$, where $H_t$ is the token representation, $w_c$ is the target-class weight, and $T$ is the number of tokens (or the number of valid tokens for masked pooling).

  \item \textbf{Method-specific combination.} The rollout signal and the head relevance signal are combined in different ways to produce several explanation modes. These are the methods currently implemented in the branch.
\end{enumerate}

The current implementation does not merely list these variants as names. Each
method is produced by the same score-construction function, so the exact
relationship between the variants is explicit. The implementation first builds
shared intermediate tensors and then chooses a different final combination rule.

\begin{enumerate}
  \item \textbf{Shared rollout construction.} The matrices passed into rollout
  are not raw attention maps from a single head. They are gradient-weighted
  layer attention matrices built in \texttt{extract\_gradient\_weighted\_attentions}.
  During the forward pass, the model returns one attention tensor per
  Transformer layer with shape \texttt{[batch, heads, tokens, tokens]}. The
  target-class logit is then backpropagated, which gives a gradient tensor of
  the same shape for each attention tensor. For each layer $l$ and head $h$,
  the implementation multiplies attention by its target-logit gradient and
  keeps only positive contributions:
  $G_{l,h} = \max(0, \nabla A_{l,h} \odot A_{l,h})$. The layer matrix used by
  rollout is the mean over heads, $A_l = \frac{1}{H}\sum_h G_{l,h}$.
  Because the attention maps are softmax probabilities and are non-negative,
  this is equivalent to $\max(0, \nabla A_{l,h}) \odot A_{l,h}$. This is the
  key methodological difference from original LeGrad, which reduces
  $\max(0, \nabla A_{l,h})$ directly without multiplying by $A_{l,h}$.

\begin{verbatim}
def build_gradient_weighted_layer_attentions(attentions, target_logit):
    retain_grad(attentions)
    target_logit.backward()

    layer_attentions = []
    for attn in attentions:
        grad = attn.grad

        weighted = grad * attn
        weighted = clamp_min(weighted, 0)

        # [batch, heads, tokens, tokens] -> [batch, tokens, tokens]
        A_l = mean(weighted, axis=heads)
        layer_attentions.append(A_l)

    return layer_attentions
\end{verbatim}

  These per-layer matrices are passed to \texttt{compute\_rollout\_attention}.
  For each layer, the implementation adds an identity matrix to preserve
  residual/self paths, row-normalizes the result, and then multiplies the
  matrices through the Transformer stack. If $I$ is the identity matrix, the
  augmented matrix is $\tilde{A}_l = \mathrm{row\_normalize}(A_l + I)$. This
  row normalization is not a softmax; it simply rescales each row to sum to
  one after the self-connection has been added. This rollout step is also not
  part of original LeGrad; it is the attention-flow component of this
  repository's hybrid explanation. The joint rollout is then
  computed by repeated batch matrix multiplication:
  $R = \tilde{A}_L \tilde{A}_{L-1} \cdots \tilde{A}_1$.

\begin{verbatim}
def compute_rollout_attention(layer_attentions):
    eye = identity(num_tokens)
    augmented = []

    for A in layer_attentions:
        A_residual = A + eye
        A_residual = A_residual / row_sum(A_residual)
        augmented.append(A_residual)

    R = augmented[0]
    for A in augmented[1:]:
        R = batch_matmul(A, R)

    return R
\end{verbatim}

  The pure rollout score is obtained by averaging over the final-token rows of
  the rollout matrix. Because HuBERT sequence classifiers do not use a ViT-like
  class token, the mean strategy treats all final tokens as possible evidence
  carriers. The averaged score is min-max scaled by
  \texttt{rollout\_to\_temporal\_relevance} and then normalized so each
  utterance has a valid distribution over time tokens.

\begin{verbatim}
raw_rollout = mean(R, axis=final_token_rows)
raw_rollout = min_max_scale(raw_rollout)
rollout_score = normalize_nonnegative(raw_rollout)
\end{verbatim}

  \item \textbf{Target-class head relevance.} The implementation captures the
  HuBERT hidden states immediately before the classification head. For the
  current SUPERB HuBERT classifier, the head is linear projection followed by
  mean pooling and a linear classifier. The code folds the projector into the
  classifier weight, so the effective class weight in HuBERT hidden-state space
  is $W_{\mathrm{eff}} = W_{\mathrm{classifier}} W_{\mathrm{projector}}$. For
  target class $c$, token $t$ receives the time-dependent logit contribution
  $H_t \cdot W_{\mathrm{eff},c} / T$, where $T$ is the number of valid tokens.
  Positive evidence is normalized first; if no positive evidence exists, the
  implementation falls back to absolute contribution, and if all contributions
  are zero it uses a uniform distribution over valid tokens.

\begin{verbatim}
def compute_head_relevance(H, logits, target_class, model, token_mask):
    W_eff = classifier.weight @ projector.weight
    w_target = W_eff[target_class]

    contribution = dot(H, w_target)        # [batch, time]
    contribution = contribution / valid_token_count(token_mask)

    positive = clamp_min(contribution, 0)
    absolute = abs(contribution)

    if sum(positive) > eps:
        return positive / sum(positive)
    if sum(absolute) > eps:
        return absolute / sum(absolute)
    return uniform_over_valid_tokens(token_mask)
\end{verbatim}

  \item \textbf{Contrastive head relevance.} The contrastive branch explains
  why the target class beats a competing class. If a contrast class is not
  supplied, the implementation chooses the highest-logit non-target class for
  each sample. It then replaces the target weight with the margin weight
  $W_{\mathrm{eff},c} - W_{\mathrm{eff},k}$ and normalizes the resulting token
  contributions with the same positive/absolute/uniform fallback.

\begin{verbatim}
def compute_contrastive_head_relevance(H, logits, target_class):
    contrast_class = argmax_non_target_logit(logits, target_class)
    W_eff = classifier.weight @ projector.weight

    margin_weight = W_eff[target_class] - W_eff[contrast_class]
    contribution = dot(H, margin_weight)
    contribution = contribution / valid_token_count(token_mask)

    return normalize_positive_abs_or_uniform(contribution)
\end{verbatim}
\end{enumerate}

\paragraph{Shared inputs used by the explanation methods.}
After the rollout and head-relevance steps above, the implementation has four
named tensors that are reused by the individual explanation methods.

\texttt{joint\_attention} is the full rollout matrix $R$ returned by
  \texttt{compute\_rollout\_attention}. It has shape
  \texttt{[batch, tokens, tokens]}. It keeps token-to-token rollout paths: an
  entry $R_{i,j}$ describes how much final token $i$ routes relevance to token
  $j$ through the stacked, gradient-weighted attention maps.

\texttt{rollout\_score} is the one-dimensional temporal rollout score derived
  from \texttt{joint\_attention}. It averages the rollout matrix over final-token
  rows, min-max scales that vector, clamps negative values, and normalizes each
  utterance so the token scores sum to one. It has shape
  \texttt{[batch, tokens]}.

\texttt{head\_relevance} is the normalized target-class classifier-head
  relevance. It is computed from the hidden states entering the head and the
  effective target-class weight $W_{\mathrm{eff},c}$. It has shape
  \texttt{[batch, tokens]} and gives the tokens that positively support the
  target class at the classifier head, with absolute/uniform fallbacks when
  positive evidence is unavailable.

\texttt{contrastive\_head} is the normalized contrastive classifier-head
  relevance. It is computed in the same space as \texttt{head\_relevance}, but
  uses the target-versus-competitor margin weight
  $W_{\mathrm{eff},c} - W_{\mathrm{eff},k}$. If the competitor is not specified,
  the highest-logit non-target class is used. It also has shape
  \texttt{[batch, tokens]}.

\begin{verbatim}
joint_attention = compute_rollout_attention(grad_attentions)

raw_rollout = rollout_to_temporal_relevance(
    joint_attention,
    strategy="mean",
)
rollout_score = normalize_token_scores(raw_rollout)

head_relevance = normalize_token_scores(
    compute_head_relevance(H, logits, target_class, model, token_mask)
)

contrastive_head, contrast_classes = compute_contrastive_head_relevance(
    H,
    logits,
    target_class,
    model,
    contrast_class,
    token_mask,
)
contrastive_head = normalize_token_scores(contrastive_head)
\end{verbatim}

\paragraph{Explanation methods.}
\begin{enumerate}
  \item \textbf{\texttt{rollout}.} This method returns the normalized rollout
  distribution directly. It is the target-agnostic baseline: it describes where
  information can flow through the Transformer, but it does not use the
  classifier head to ask which tokens support a particular emotion class.

\begin{verbatim}
scores["rollout"] = rollout_score
\end{verbatim}

  \item \textbf{\texttt{level3}.} This method keeps the original product-style
  Level 3 formulation used in the branch. It intersects the target-agnostic
  rollout map with target-class evidence from the classifier head. A token is
  therefore important only when it is both reachable through the rollout graph
  and positively contributes to the target class at the head.

\begin{verbatim}
scores["level3"] = normalize_nonnegative(
    rollout_score * head_relevance
)
\end{verbatim}

  \item \textbf{\texttt{level3-contrastive}.} This variant uses the same
  product rule, but replaces target-class head relevance with target-versus-
  competitor margin relevance. It highlights tokens that support the target
  class relative to the selected competing class, rather than tokens that merely
  support the target logit in isolation.

\begin{verbatim}
scores["level3-contrastive"] = normalize_nonnegative(
    rollout_score * contrastive_head
)
\end{verbatim}

  \item \textbf{\texttt{head-conditioned-rollout}.} This method does not average
  every row of the joint rollout matrix. Instead, it uses the normalized
  target-class head relevance as a seed distribution over final tokens and
  propagates that seed backward through the joint rollout matrix. If
  $R_{i,j}$ describes the rollout path from final token $i$ to input token $j$,
  the conditioned score for token $j$ is
  $\sum_i s_i R_{i,j}$, where $s_i$ is the normalized head-relevance seed.

\begin{verbatim}
def rollout_from_token_relevance(R, token_relevance):
    seeds = clamp_min(token_relevance, 0)
    seeds = normalize_or_uniform(seeds)

    conditioned = batch_matmul(seeds[:, None, :], R)
    conditioned = squeeze(conditioned, axis=1)
    conditioned = clamp_min(conditioned, 0)

    return conditioned / sum(conditioned)

scores["head-conditioned-rollout"] = rollout_from_token_relevance(
    R,
    head_relevance,
)
\end{verbatim}

  \item \textbf{\texttt{contrastive-conditioned-rollout}.} This is the
  contrastive version of the conditioned rollout method. The seed distribution
  is the target-versus-competitor head relevance, and the rollout matrix
  propagates those contrastive seeds through the Transformer attention paths.

\begin{verbatim}
scores["contrastive-conditioned-rollout"] = rollout_from_token_relevance(
    R,
    contrastive_head,
)
\end{verbatim}
\end{enumerate}

The implementation is therefore a lightweight, auditable hybrid of
attention-gradient weighting, rollout, and classifier-head relevance. It is
useful for experiments because it is easy to compute, easy to compare across
methods, and directly produces a temporal relevance curve over the input
waveform. It should not be reported as an exact reproduction of the LeGrad
algorithm.

\section{Notebook-by-Notebook Technical Overview}

\subsection{Notebook 00 --- Model Used and Classification Head Architecture}

\subsubsection*{Purpose}
Document the active model architecture used by the current transformer-relevance
pipeline and explain why the classification head design matters for a faithful
LeGrad adaptation.

\subsubsection*{Technical idea}
The notebook loads the model configuration for
\texttt{superb/hubert-base-superb-er}, displays the relevant architecture
fields, renders the HuBERT sequence-classification path, and prints the
installed \texttt{transformers} forward-code excerpt where
\texttt{use\_weighted\_layer\_sum}, \texttt{layer\_weights},
\texttt{projector}, temporal pooling, and \texttt{classifier} appear. It then
compares this design with the ViT/CLIP path used by the official LeGrad code.

\subsubsection*{Decisions made}
\begin{itemize}
  \item The notebook treats \texttt{superb/hubert-base-superb-er} as the active
  model for the current relevance scripts, while displaying the older YAML model
  source for traceability.
  \item The notebook inspects the real checkpoint tensor file to report learned
  \texttt{layer\_weights} and the concrete projector/classifier parameter
  shapes without relying only on configuration metadata.
  \item The notebook separates the convolutional frontend issue from the
  classification-head issue: the frontend affects token-to-time mapping, while
  the head affects how a layer-wise LeGrad score should be defined.
\end{itemize}

\subsubsection*{Outcome}
Notebook \path{notebooks/00_model_used.ipynb} was executed and now provides a
visual explanation of the active model. The run confirmed a
\texttt{HubertForSequenceClassification} configuration with 12 Transformer
layers, 12 attention heads, hidden size 768, classifier projection size 256,
four emotion labels, and \texttt{use\_weighted\_layer\_sum=True}. It also
prints the installed forward-code excerpt showing the learned hidden-layer
mixture before projection, pooling, and classification. The checkpoint
inspection confirms that the learned mixture uses 13 entries, one for $H_0$ and
one for each of the 12 Transformer outputs, before the 768-to-256 projector and
4-class emotion classifier.

\subsection{Notebook 01 --- SpeechXAI Single Example}

\subsubsection*{Purpose}
Establish a complete SpeechXAI explanation and deletion-evaluation workflow on one IEMOCAP utterance before integrating the LeGrad-inspired method.

\subsubsection*{Technical idea}
SpeechXAI uses WhisperX to transcribe the utterance and align each word with its actual start and end time. The same timestamp objects are passed into SpeechXAI's leave-one-out explainer, ensuring that each attribution score refers to the word interval shown in the visualization. SpeechXAI measures importance as the confidence change caused by replacing a word region with silence. The paper describes zeroing the samples inside the WhisperX word interval. The released repository silences a larger interval: it starts 100\,ms before the WhisperX word start and ends 40\,ms after the WhisperX word end. Thus, the attribution assigned to one word is calculated after silencing the word together with these adjacent audio samples.

\subsubsection*{Decisions made}
\begin{itemize}
  \item A pretrained wav2vec 2.0 emotion classifier is used so the project can focus on explanation quality rather than classifier training.
  \item One utterance is used to keep this notebook focused on the explanation workflow and visualization.
  \item The original predicted class is the attribution and deletion target.
  \item Silence replacement preserves the waveform length.
  \item SpeechXAI's WhisperX transcription is the only source of word text and timing; no character-based or uniform timing approximation is used.
  \item Transcription is performed once, and the returned word segments are reused by the explainer and plot so words, timestamps, and scores remain index-aligned.
  \item WhisperX's per-word transcription confidence is retained in the attribution table so uncertain alignments remain visible during interpretation.
  \item The paper's Section 2.1 and Appendix B.1 describe zeroing word segments derived from WhisperX timestamps to preserve the original recording length. They do not specify temporal padding around those segments.
  \item The released repository's \texttt{utils\_removal.py} defines \texttt{a=100} and \texttt{b=40}. For a WhisperX word interval $[t_{\mathrm{start}},t_{\mathrm{end}}]$, \texttt{remove\_word} and \texttt{remove\_specified\_words} silence $[t_{\mathrm{start}}-100\,\mathrm{ms},t_{\mathrm{end}}+40\,\mathrm{ms}]$.
  \item The notebook reproduces the repository implementation rather than silently replacing it with the paper-level interval description.
  \item The visualization distinguishes two regions: the lexical interval reported by WhisperX and the larger interval actually silenced by the repository.
  \item A separate overlap diagnostic displays every effective mask on its own row and computes sample-level mask depth. Regions covered by two or more masks are highlighted explicitly.
  \item Attribution bars use SpeechXAI's raw confidence-drop values. They are not rescaled to make the largest word equal to one, because that would hide the magnitude of the model response.
  \item The waveform-loading path used by the project is checked against SpeechXAI's pydub-loading path. Both must produce the same predicted class and probability vector before the explanation is accepted.
  \item Word intervals must be chronological, have positive duration, remain inside the audio, and correspond one-to-one with SpeechXAI's returned features and scores.
  \item The notebook rejects an utterance when its first aligned word begins before SpeechXAI's 100\,ms pre-word padding. The upstream removal function does not clamp this slice at zero, so accepting such an example could make the displayed mask differ from the perturbation actually evaluated.
  \item The final evaluation uses the top-$k$ SpeechXAI words to define the shared masking duration.
  \item Overlapping padded masks are counted by their union duration so the audio-removal budget is not inflated.
  \item The notebook links the upstream implementation modules beside the attribution code. The LOO explainer calls single-word removal, the removal utility performs waveform replacement, and the faithfulness evaluator invokes multi-word removal. The linked files are \path{loo_speech_explainer.py}, \path{utils_removal.py}, and \path{faithfulness_measures_speech.py}.
  \item Notebook 01 writes no result files; all results remain visible in the executed notebook.
\end{itemize}

\subsubsection*{What was tested}
The selected audio was loaded and classified, then transcribed and aligned with SpeechXAI's WhisperX function. The returned segments were passed directly into SpeechXAI's leave-one-out explainer and checked against the returned feature order. Timestamp validity, word-score alignment, and the 100\,ms first-word safety condition were validated. The original prediction obtained from SpeechXAI's pydub representation matched the prediction obtained from the waveform shown in the notebook. The first visualization used real word locations, effective silence masks, raw confidence-drop attribution values, and per-word transcription confidence. The second visualization inspected padded removal intervals and computed overlap depth at waveform-sample resolution. Top-$k$ SpeechXAI deletion was compared with repeated random masking using the same union duration.

\subsubsection*{Outcome}
The notebook executed successfully for utterance \texttt{Ses01F\_impro01\_F005}. WhisperX located the first word at 0.594 seconds, preserving the initial silence and later pauses. The returned word-confidence values ranged from 0.760 to 0.980. The classifier predicted anger with confidence 0.8781. SpeechXAI assigned the largest raw confidence drop, 0.4218, to ``what's''. The overlap diagnostic found six distinct regions covered by multiple padded masks, with a maximum depth of two simultaneous masks. The visualizations, attribution table, consistency checks, and shared-duration evaluation all completed successfully.

\subsubsection*{Remaining issues}
The current validation covers one utterance only. WhisperX confidence is useful diagnostic information but is not a guarantee that every boundary is correct. The 100/40\,ms expansion is present in the released Python code but is not documented in the paper, so future experiments should state explicitly whether they reproduce the repository behavior or use unpadded paper-level word intervals. Utterances with a word beginning inside the first 100\,ms are currently rejected because of the upstream SpeechXAI slicing behavior; supporting them safely would require clamping the library's perturbation boundary. The pipeline must be run on a controlled multi-audio set before drawing comparative conclusions.

\subsection{Notebook 02 --- Attention Access Test}

\subsubsection*{Purpose}
Determine whether the wav2vec 2.0 Transformer attention maps required by the LeGrad-inspired method are accessible.

\subsubsection*{Technical idea}
The active classifier is a Hugging Face \texttt{Wav2Vec2ForSequenceClassification} model. The diagnostic step requests attention outputs directly from its Transformer encoder and inspects their dimensions and numerical behavior.

\subsubsection*{Decisions made}
\begin{itemize}
  \item Attention access is treated as a required checkpoint before developing LeGrad.
  \item The Hugging Face classifier is accessed directly, matching the model-loading path used by SpeechXAI.
  \item Attention tensors should be inspected across layers and heads and visualized for at least one layer-head pair.
\end{itemize}

\subsubsection*{What was tested}
The configured SUPERB classifier was run on the same IEMOCAP utterance used in Notebook 01 with attention output enabled. The notebook inspected every returned layer, validated finite values, printed tensor dimensions, and visualized one final-layer attention head.

\subsubsection*{Outcome}
Attention access succeeded. The model returned 12 attention tensors, one per Transformer layer, each shaped \texttt{[1, 12, 201, 201]} for the selected utterance.

\subsubsection*{Remaining issues}
Notebook 03 must confirm that these returned attention tensors retain usable, nonzero gradients for the original predicted-class logit.

\subsection{Notebook 03 --- Attention Gradient Test}

\subsubsection*{Purpose}
Verify that gradients of a target class score with respect to Transformer attention maps can be obtained.

\subsubsection*{Technical idea}
LeGrad uses sensitivity to attention maps as its explanation signal. For speech, the target quantity is the classifier score for the original predicted emotion, and the required diagnostic is $\partial s_{\hat{y}}/\partial A_l$ for one or more Transformer layers. The repository now keeps two families separate: older hybrid rollout baselines multiply positive gradients by attention values before rollout, while the LeGrad-HuBERT modes in Notebook 04 use positive attention gradients directly as the relevance signal.

\subsubsection*{Decisions made}
\begin{itemize}
  \item The target is the original predicted-class logit or score.
  \item Returned attention tensors will retain gradients during backpropagation.
  \item Gradient shape, magnitude, finite values, and nonzero proportion will be checked before attribution aggregation.
  \item Logits are preferred over probabilities when probability saturation weakens gradients.
\end{itemize}

\subsubsection*{What was tested}
The notebook scaffold and reusable gradient-retention utilities are present.

\subsubsection*{Outcome}
The conceptual gradient path is defined. A real model forward and backward pass remains to be validated after Notebook 02 succeeds.

\subsubsection*{Remaining issues}
Gradient availability depends on how the wav2vec 2.0 attention implementation returns intermediate tensors. Hooks may be required if returned attentions are detached from the computational graph.

\subsection{Notebook 04 --- LeGrad-Inspired Temporal Explanation}

\subsubsection*{Purpose}
Implement and validate explicit LeGrad-style HuBERT temporal relevance modes
on one real IEMOCAP utterance.

\subsubsection*{Technical idea}
Notebook 04 now uses the implementation in
\path{src/explainers/transformer_relevance/legrad_hubert.py}. The LeGrad-HuBERT
modes compute the positive attention-gradient signal directly,
\[
\max(0, \partial s / \partial A_l),
\]
then reduce that signal over attention heads and temporal query positions before
aggregating layers. Two score definitions are exposed: a final-score variant using the real HuBERT score
\[
s_c = \mathrm{classifier}(\mathrm{pool}(\mathrm{projector}(H_{\mathrm{mix}})))_c
\]
and a layer-local variant using
\[
s_{l,c} = \mathrm{classifier}(\mathrm{pool}(\mathrm{projector}(H_l)))_c.
\]
For both score definitions, relevance is reduced over heads and query tokens so
that source/key temporal tokens remain as the explanation axis.

\subsubsection*{Decisions made}
\begin{itemize}
  \item Four long, explicit mode names were added so result files identify the
  score source, formula, layer aggregation, and token axis.
  \item The implemented formula for LeGrad-HuBERT modes is
  \texttt{ReLU(d score / d attention)}, not
  \texttt{ReLU((d score / d attention) * attention)}.
  \item Both mean layer aggregation and renormalized HuBERT-layer-weighted
  aggregation are implemented, so the analysis can choose between a
  LeGrad-like mean and a HuBERT-aware weighting.
  \item Each mode returns sum-normalized temporal scores for quantitative
  evaluation and min-max-normalized visual scores for paper-style plots.
  \item The primary notebook visualization uses the layer-local score with
  renormalized HuBERT layer weights, while the notebook also reports the
  final-score variants.
  \item Temporal outputs use small fixed-duration bins or token intervals that can participate in duration-matched deletion evaluation.
\end{itemize}

\subsubsection*{What was tested}
Notebook 04 was executed on utterance \texttt{Ses01F\_impro01\_F005}. It loaded
the HuBERT classifier, computed all four LeGrad-HuBERT modes, displayed the
predicted-class logits, printed intermediate layer-relevance shapes, displayed
the renormalized HuBERT layer weights and layer-local target logits, validated
finite sum-normalized evaluation scores and min-max visual scores, listed the
top temporal tokens, and rendered both a Notebook-01-style waveform relevance
plot and the existing spectrogram timeline plot.

\subsubsection*{Outcome}
The notebook completed successfully. The selected utterance produced 12
Transformer attention layers and 201 temporal tokens. Both
\texttt{final\_score\_layer\_relevance} and
\texttt{layer\_local\_score\_layer\_relevance} had shape
\texttt{(1, 12, 201)}. The renormalized HuBERT Transformer-layer weights summed
to 1.0, and each final temporal relevance vector summed to 1 within numerical
tolerance.

\subsubsection*{Remaining issues}
The current token-to-time mapping remains the existing uniform approximation
over the waveform duration. A later refinement could map HuBERT feature frames
through the convolutional frontend stride and receptive field more explicitly.

\subsection{Notebook 05 --- Quantitative Evaluation}

\subsubsection*{Purpose}
Run the complete multi-audio comparison in one coherent experiment. This notebook owns batch explanation generation, common interval representation, equal-duration masking, deletion metrics, and aggregate quantitative plots.

\subsubsection*{Technical idea}
The validated SpeechXAI workflow from Notebook 01 and LeGrad-inspired workflow from Notebook 04 are applied to the same controlled IEMOCAP set. SpeechXAI top-$k$ words define a union duration $X$. SpeechXAI masks, LeGrad-inspired top bins, and repeated random bins are then compared after silencing exactly $X$ seconds while preserving waveform length.

\subsubsection*{Decisions made}
\begin{itemize}
  \item Batch SpeechXAI generation is part of evaluation rather than a separate notebook because its single-audio behavior is already established in Notebook 01.
  \item Interval unification and masking are part of the evaluation operation rather than independent experimental outcomes.
  \item Both explanation methods use the common columns \texttt{audio\_id}, \texttt{start}, \texttt{end}, and \texttt{score}.
  \item SpeechXAI is the source of duration $X$ because its native unit is a variable-duration word.
  \item Top-$k$ values are $1$, $2$, $3$, and $5$.
  \item LeGrad-inspired bins are selected by score until their union duration reaches $X$.
  \item Random bins use the same duration and are repeated 20 times.
  \item The main metric is absolute confidence drop for the original predicted class.
  \item Relative confidence drop and prediction flip rate are secondary metrics.
  \item Aggregate tables, deletion curves, boxplots, duration checks, and flip rates remain in this notebook because they are direct outputs of the quantitative experiment.
\end{itemize}

\subsubsection*{What was tested}
Notebook 01 validated real SpeechXAI attribution, silence masking, union-duration accounting, confidence drop, and repeated same-duration random masking on one utterance. The consolidated Notebook 05 currently contains the protocol scaffold that will extend those validated operations to SpeechXAI and LeGrad across the evaluation set.

\subsubsection*{Outcome}
The evaluation schema and fairness constraints are defined and validated at single-audio scale. The previous standalone SpeechXAI, interval-masking, deletion, and final-summary placeholders have been removed because their responsibilities now belong to this one quantitative experiment.

\subsubsection*{Remaining issues}
The complete run requires a working LeGrad example from Notebook 04, a documented evaluation set, and execution across all selected audios.

\subsection{Notebook 06 --- Qualitative Examples}

\subsubsection*{Purpose}
Create interpretable examples that show where SpeechXAI and LeGrad agree or disagree and how masking changes model confidence.

\subsubsection*{Technical idea}
Each example combines the waveform, word-level SpeechXAI intervals, the LeGrad temporal relevance curve, selected masking regions, and before/after classifier confidence.

\subsubsection*{Decisions made}
\begin{itemize}
  \item Examples will be selected from quantitative outcomes rather than chosen only for visual appeal.
  \item The final set should include agreement, LeGrad-favored, SpeechXAI-favored, and difficult cases.
  \item Interpretations will distinguish model faithfulness from human plausibility.
\end{itemize}

\subsubsection*{What was tested}
The notebook structure and plotting requirements have been defined.

\subsubsection*{Outcome}
The qualitative analysis plan is ready.

\subsubsection*{Remaining issues}
Representative cases can be selected after Notebook 05 produces the complete deletion results. Aggregate tables remain in Notebook 05; this notebook is reserved for case-level interpretation.

\section{Evaluation Design}

The final quantitative evaluation follows these steps:

\begin{enumerate}
  \item SpeechXAI returns importance scores for word-aligned audio segments.
  \item For each audio and each $k \in \{1,2,3,5\}$, the top-$k$ SpeechXAI words are selected.
  \item The selected words define a total duration $X$ seconds.
  \item Those SpeechXAI word intervals are replaced by silence, and the confidence drop for the original predicted class is measured.
  \item LeGrad-inspired time bins are ranked by relevance, and the highest-ranked bins whose total duration equals $X$ are selected.
  \item Those LeGrad-inspired regions are replaced by silence, and the same confidence drop is measured.
  \item Random time bins with total duration $X$ are selected repeatedly to form a random deletion by silence masking baseline.
  \item All methods use the same audio, classifier, target class, masking operation, and removed duration.
\end{enumerate}

SpeechXAI naturally explains words, while the LeGrad-inspired method explains Transformer time bins or tokens. Equalizing the total occluded duration makes the comparison focus on where each method assigns relevance.

The main metric is:

\[
\mathrm{drop}
= p_{\mathrm{original}}(\hat{y})
- p_{\mathrm{masked}}(\hat{y}),
\]

where $\hat{y}$ is the class predicted from the original audio. Secondary metrics are relative confidence drop and prediction flip rate. Mean masked duration is reported for every top-$k$ setting.

\section{Methodological Decisions Summary}

\begin{longtable}{p{0.20\textwidth} p{0.34\textwidth} p{0.38\textwidth}}
\toprule
\textbf{Component} & \textbf{Decision} & \textbf{Reason} \\
\midrule
Dataset & IEMOCAP speech, transcripts, and emotion labels & It is a standard speech emotion resource and provides the audio and annotations required by the selected classifier. \\

Classifier & SUPERB \texttt{wav2vec2-base-superb-er} checkpoint used by SpeechXAI & Using the same classifier removes model choice as a confound when comparing SpeechXAI and LeGrad-inspired explanations. \\

Model label space & Anger, happiness/excitement, neutral, sadness & These are the four outputs supported by the selected pretrained classifier. \\

Target class & Original predicted class & Faithfulness is evaluated relative to the model decision being explained. \\

SpeechXAI unit & Word-aligned audio segments & This is SpeechXAI's native, human-readable explanation unit. \\

LeGrad unit & Transformer tokens or fixed temporal bins & Attention-gradient relevance is naturally produced at the encoder-token level. \\

LeGrad-HuBERT formulation & Positive attention gradients; final-score and layer-local score variants are both implemented & This is the explicit HuBERT adaptation of the LeGrad reduction, while the older gradient-attention rollout modes remain available as baselines. \\

Shared duration & $X$ seconds from SpeechXAI top-$k$ words & This makes word-level and time-bin explanations comparable under equal audio removal. \\

Top-$k$ values & $1$, $2$, $3$, and $5$ words & These settings provide a compact deletion curve while respecting SpeechXAI's native output. \\

Masking method & Silence replacement with unchanged waveform length & It provides a simple perturbation and avoids duration changes that could independently affect the classifier. \\

Random baseline & Same-duration random bins, repeated 20 times & Repetition estimates baseline variance and provides a stable reference. \\

Main metric & Confidence drop for the original predicted class & It measures how strongly the selected regions support the model's original decision. \\

Secondary metrics & Relative confidence drop and prediction flip rate & They complement absolute confidence drop across samples with different confidence levels. \\

Notebook 01 scope & One utterance, in-memory outputs & It establishes an understandable baseline and avoids mixing pipeline validation with dataset-scale evaluation. \\
\bottomrule
\end{longtable}

\end{document}
